En este capítulo se recoge el diseño de la arquitectura del sistema informático, la parametrización del software base (opcional), el diseño físico de datos y el diseño detallado de la interfaz de usuario.

\section{Arquitectura del Sistema}
\label{sec:SystemArchitecture}
\input{./sections/development/design/systemarchitecture.tex}

\section{Parametrización del software base}
En esta sección (opcional), se detallan las modificaciones a realizar sobre el software base, que son requeridas para la correcta construcción del sistema. En esta sección incluiremos las  actuaciones necesarias sobre la interfaz de administración del sistema, sobre el código fuente o sobre el modelo de datos.

\section{Diseño Físico de Datos}
En esta sección se define la estructura física de datos que utilizará el sistema, a partir del modelo de conceptual de clases, de manera que teniendo presente los requisitos establecidos para el sistema de información y las particularidades del entorno tecnológico, se consiga un acceso eficiente de los datos. La estructura física se compone de tablas, índices, procedimientos almacenados, secuencias y otros elementos dependientes del SGBD a utilizar.

\section{Diseño de la Interfaz de Usuario} 
En esta sección se detallarán los principios a tener en cuenta a la hora de implementar la interfaz de usuario del sistema, incluyendo el conjunto de estilos (colores, tipografías, iconografía, etc.) y el conjunto de componentes visuales que se utilizarán. Se incluirá, además, un prototipo de alta fidelidad con el diseño de algunas de las pantallas del sistema.
